% ------------------------------------------------------------------------------
% attointern: new version of attointer
% ------------------------------------------------------------------------------
\documentclass[12pt]{article}
\usepackage[a4paper,left=2.5cm,right=2.5cm,top=2.5cm,bottom=2cm]{geometry}
\usepackage[T1]{fontenc}
\usepackage[latin1]{inputenc}
\usepackage{newtxtext}
\usepackage[ttscale=0.875]{libertine}
\usepackage[libertine,cmintegrals,cmbraces]{newtxmath}
\usepackage{authblk}
\usepackage{graphicx}
\usepackage{setspace}
\usepackage[round]{natbib}
\usepackage[table,xcdraw,svgnames]{xcolor}
\usepackage{booktabs}
\usepackage{tabu}
\usepackage{soul}
\usepackage{url}
\usepackage{lineno}


\DeclareMathOperator{\logmod}{logmod}%
\DeclareMathOperator{\expmod}{expmod}%
\DeclareMathOperator{\sign}{sign}%      2016-12-16T15:04:44


\newcommand{\sub}[1]{\textsubscript{#1}}
\newcommand{\super}[1]{\textsuperscript{#1}}
\newcolumntype{C}[1]{>{\centering\let\newline\\\arraybackslash\hspace{0pt}}m{#1}}

\setstcolor{red}
%~ \newcommand{\correc}[2]{{\color{red} #1}\,{\color{blue} #2}}
\newcommand{\correc}[2]{{\color{red}\st{#1}}\,{\color{blue} #2}}
%\newcommand{\correc}[1]{{\color{blue} #1}}

% -----------------------------------------------------------------------------
% degree
% -----------------------------------------------------------------------------
\newcommand{\degree}{\ensuremath{^\circ}}
\newcommand{\ape}{~}
\newcommand{\pretit}[1]{}
\newcommand{\retrieved}{Retrieved  in}
\newcommand{\available}{Available  at}
\newcommand{\submitted}{submitted}
\newcommand{\edone}{\rst }
\newcommand{\edtwo}{\nd\ }
\newcommand{\edthree}{\rd }
\newcommand{\ednth}{\th }



\clubpenalty=10000
\widowpenalty=10000


\begin{document}

\setstretch{1.6}
\linenumbers




\title{Path shadowing correction}

  \author[1]{Nelson Lu�s Dias}

  

\affil[1]{Department of Environmental Engineering, Federal University of Paran�,  Brazil. }


\maketitle












{ \linelabel{lin:begin-explainshadow} When shadowing corrections are applied
  \linelabel{lin:only} (only to the CSAT3s) to data that passed QC , this is always done
  \emph{before} any other processing: in these cases, other procedures
  such as calculation of planar fit rotation matrices and the
  rotations themselves, calculation of means and statistics, etc., are
  performed on the corrected data.  \linelabel{lin:end-explainshadow}}

From the $(u,v,w)$ in an orthogonal system measured by the CSAT3,
we recover the path components
\cite[Eq. (10)]{horst.oncley--correctionstoinertial} from
\begin{align}
  u_a &= -u \cos(\phi) + w\sin(\phi), \\
  u_b &= \frac{u + \sqrt{3}v}{2}\cos(\phi) + w\sin(\phi), \\
  u_c &= \frac{u - \sqrt{3}v}{2}\cos(\phi) + w\sin(\phi),
\end{align}
where $\phi = \pi/3$ is the angle between the CSAT3 paths and the horizontal.
Next, we calculate the instantaneous absolute value of the (uncorrected) wind
vector,
\begin{equation}
S = \sqrt{u^2 + v^2 + w^2},\label{eq:Sfrom-uvw}
\end{equation}
and then, for each path component, the angle between the 
wind vector and that path:
\begin{equation}
\theta_{\alpha} = \arccos(u_{\alpha}/S),
\end{equation}
where $\alpha = a, b, c$. Then the path components are corrected by \citep{horst.semmer.ea--correction}
\begin{equation}
u_{\alpha} := \frac{u_{\alpha}}{0.84 + 0.16\sin(\theta_{\alpha})}.
\end{equation}
where $:=$ stands for the usual assignment in computer languages.
The orthogonal components are then recovered by an inverse rotation
\cite[Eq. (9)]{horst.oncley--correctionstoinertial}:
\begin{align}
  u &= \frac{-2u_a + u_b + u_c}{3\cos(\phi)}, \label{eq:new-u}\\
  v &= \frac{u_b - u_c}{\sqrt{3}\cos(\phi)},\label{eq:new-v}\\
  w &= \frac{u_a + u_b + u_c}{3\sin(\phi)}.\label{eq:new-w}
\end{align}
Note  that the newly recovered $(u,v,w)$ is different from the
raw measurement; therefore, one needs to iterate trough
(\ref{eq:Sfrom-uvw})--(\ref{eq:new-w}) until convergence is achieved
in all components. In this work, the convergence criterion is that the
maximum difference over the three components is less than
$\mathrm{0.001\,m\,s^{-1}}$, which is well below the stated accuracy
($0.01\,\mathrm{m\,s^{-1}}$) of the CSAT3.

{ \linelabel{lin:begin-refshadow} Campbell's CSAT3B manual
  \cite[][section 7.1.3]{csi--csat3b} mentions that the option to
  apply the shadowing correction uses 3
  iterations. \cite{zhou2018recovery} give a very detailed description
  of the shadowing correction algorithm, and impose, as done here, a
  convergence criterion of $\mathrm{1\,mm\,s^{-1}}$ on each velocity
  component. They report a maximum of 3 iterations until
  covergence. In our case we found that a maximum of 6 iterations were
  needed. \linelabel{lin:end-refshadow}}


\bibliographystyle{plainnat-lastname}
\bibliography{abbrevs,all}

\end{document}
